\documentclass[conference]{IEEEtran}
\usepackage[utf8]{inputenc}
\usepackage[T1]{fontenc}
\usepackage{cite}
\usepackage{amsmath,amssymb}
\usepackage{graphicx}
\usepackage{booktabs}
\usepackage{url}

\title{Does Deterministic Generate--Validate--Repair Reduce Unsafe LLM‑Generated Network Changes? A Confirmatory Paired Study}

\author{
\IEEEauthorblockN{Autonomous AI Researcher}
\IEEEauthorblockA{CNSM 2026 Agentic AI Researcher Track}
}

\begin{document}

\maketitle

\begin{abstract}
We preregistered and executed a sealed paired confirmatory study (n = 200 intent\ensuremath{\rightarrow}configuration tasks) to test whether a deterministic generate--validate--repair (GVR) pipeline that consults a deterministic NetOps validator reduces validator‑detected unsafe or semantically incorrect network changes relative to an LLM‑only generation pipeline. The run used a single hosted model (gpt‑5‑mini) with adapter‑enforced shared‑initial generation semantics and a primary deterministic validator (deterministic\_netops\_validator\_v5). Execution completed all planned pairs (completed\_pairs = 200) consuming 200 model calls (one shared initial generation per task); no repair calls were issued because every initial candidate passed the validator. Paired contingency: n11 = 200, n10 = 0, n01 = 0, n00 = 0; estimate = 0.0; exact McNemar two‑sided p = 1.0; 95\% bootstrap percentile CI = [0.0, 0.0] (10,000 resamples, seed = 1729). Because the baseline validator‑detected unsafe incidence was 0\% on this task bank, the preregistered >=50\% relative‑reduction hypothesis could not be meaningfully evaluated. We report the sealed design, execution accounting (start: 2026‑08‑30T11:50:34.530059Z; end: 2026‑08‑30T12:12:07.541650Z), the Mininet 10\% holdout (20 tasks) that produced zero validator--emulator disagreements, and an artifact‑grounded discussion of operational implications and threats to validity (preregistration id: cnsms2026‑gvr‑effectiveness‑02).
\end{abstract}

\section{Introduction}
Automating intent\ensuremath{\rightarrow}configuration translation with generative language models can increase NetOps productivity but risks producing incorrect or unsafe changes that cause outages or policy violations. Prior work therefore proposes grounding, deterministic validation, and repair loops (generate--validate--repair, GVR) to detect and correct unsafe candidates before application. The operational benefit of automated repair is conditional on three factors: (i) baseline prevalence of validator‑detectable failures from the chosen generator/prompting policy; (ii) validator coverage and fidelity; and (iii) the available repair budget and repair‑prompt effectiveness. We preregistered a narrowly scoped confirmatory question: Does an automated, deterministic GVR pipeline that consults a deterministic NetOps validator materially reduce validator‑detected unsafe or semantically incorrect network changes relative to an LLM‑only generation pipeline on a curated 200‑task intent\ensuremath{\rightarrow}configuration bank? The study was designed as a paired experiment that fixes the generator output per task and isolates the marginal effect of invoking a validator+repair on the same candidate.

\section{Related Work}
LLM failure modes, mitigation primitives, and NetOps prototypes. Systematic surveys and empirical studies across LLM applications have repeatedly documented hallucinations and factual errors, substantial prompt sensitivity, and non‑determinism in generated outputs; these phenomena motivate a family of mitigations including retrieval grounding, verifier‑assisted repair, and constrained or guarded agent architectures [10.1007/s11704-026-60308-3, https://openalex.org/W4399803256, https://openalex.org/W4402860127]. In the NetOps domain specifically, multiple prototype systems demonstrate that LLMs can translate natural‑language intent into configuration fragments or workflows, but these works commonly stop short of instrumented, production‑style failure‑rate studies that quantify end‑to‑end operational risk [https://openalex.org/W4399629579, https://openalex.org/W4403635865].

Generate--Validate--Repair (GVR) and deterministic validators. Architectural proposals and early implementations advocate integrating LLM generation with deterministic, rule‑based validators (Batfish‑style or intent/constraint checkers) to prevent straightforward syntactic and semantic violations before application; GVR loops then attempt automated repairs when validators report violations [https://openalex.org/W4410125989, doi:10.36227/techrxiv.177004277.71795055/v1]. These integrations aim to convert an open‑ended generative problem into a closed‑loop verification problem where correctness is operationalized by validator predicates (for example: required\_command\_count, parsed\_command\_count, violation\_codes). The literature also notes that validator coverage is crucial: deterministic checkers can be strong for the class of constraints they encode (topology‑level invariants, required sequences), but may omit device‑level idiosyncrasies or data‑plane effects only visible under emulation or on real hardware [https://openalex.org/W4399629579]. That gap motivates the use of secondary end‑to‑end emulators or sampled real‑device checks as contamination detectors in empirical evaluations.

Agentic workflows, tool‑use brittleness, and red‑team findings. Studies of agentic LLM systems and tool‑using agents highlight new failure modes: tool‑description or tool‑registry injection, cross‑tool ambiguity, and persona/history interactions that change how agents interpret verifier responses or construct repairs [10.21203/rs.3.rs-10646209/v1, 10.2139/ssrn.7203631]. Red‑team and audit reports show that apparently secure tool‑assisted pipelines can be manipulated by adversarially crafted tool definitions or chained actions, producing dangerous configuration changes even when individual components seem benign. These findings imply that GVR effectiveness depends not only on validator logic but on the larger tool and agent context in which repair prompts and validator outputs are framed.

Evaluation gaps and the external‑validation problem. Recent meta‑analyses and reviews emphasize a recurring external‑evaluation gap: many NetOps prototypes and LLM evaluations use synthetic or narrowly curated tasks and do not measure how validator‑detected safety maps to real end‑to‑end risk across vendor devices, multi‑vendor behaviors, or adversarial inputs [https://openalex.org/W4411523078, https://openalex.org/W4411735779]. That gap motivates three experimental priorities: (1) instrumented paired studies that fix generator sampling to isolate validator/repair effects, (2) inclusion of emulator or real‑device holdouts to detect validator blind spots, and (3) reporting of repair budgets and model‑call costs so practitioners can assess operational tradeoffs.

Prompt sensitivity, sampling semantics, and experimental design implications. Empirical work on prompt sensitivity and generation non‑determinism shows that sampling policy (temperature, number of samples) and prompt templates materially affect observed correctness rates and the opportunities available to a repair loop [https://openalex.org/W4402860127]. Consequently, a paired shared‑initial design---where the same initial candidate is reused across baseline and guarded arms---removes sampling variance and isolates the marginal effect of validation+repair on a fixed generator output. This design choice trades generality for controlled inference and is appropriate when the research question targets the marginal utility of repair for fixed generator behavior.

Synthesis and positioning. The prior work corpus collectively supports GVR as a promising architectural pattern but leaves unresolved empirical questions about when deterministic repair provides measurable operational benefit versus when efforts should focus on improving validator fidelity, diversifying generator sampling, or hardening agents against tool‑description attacks. Our preregistered paired study was designed to address this specific empirical gap by fixing generator outputs (shared‑initial semantics), enforcing a single‑repair budget, and instrumenting a randomized emulator holdout as a contamination detector; the study therefore complements prior prototype and audit work by producing an externally auditable, artifact‑grounded datum about GVR's marginal effect under a constrained but clearly specified operational regime.

\section{Methodology}
Preregistration and inferential plan. We froze the experimental plan under preregistration id cnsms2026‑gvr‑effectiveness‑02 before any model calls. The confirmatory design was paired and deterministic: n = 200 tasks (four topology clusters \ensuremath{\times} 50 tasks), one shared deterministic initial LLM generation per task, and a repair budget of at most one automated repair call per task. The primary estimand was the paired\_success\_rate\_difference\_guarded\_minus\_baseline and the prespecified test was an exact two‑sided McNemar test (\ensuremath{\alpha} = 0.05). Confidence intervals for the paired difference were computed by the bootstrap percentile method with 10,000 resamples and bootstrap\_seed = 1729. The analysis executor was paired\_binary\_analysis\_v1 as recorded in the analysis artifacts.

Adapter, model, and sampling semantics. The run used a single hosted model declared as gpt‑5‑mini (provider label: openai\_responses). Model configuration metadata (execution/model\_configuration.json) records requested\_model = "gpt-5-mini", max\_output\_tokens = 2000, maximum\_attempts\_per\_call = 1, and temperature = null (deterministic sampling enforced by the adapter). Crucially, the hosted\_netops\_gvr\_v1 adapter enforced shared\_initial generation semantics: for each task the system generated exactly one initial candidate which was deterministically reused for both the baseline (LLM‑only) and guarded (GVR) conditions. This shared‑initial design eliminates per‑condition sampling variance and isolates the incremental effect of invoking validation and an allowed repair step on the identical generator output. Adapter accounting in deterministic\_reconciliation.json documents hosted\_model\_call\_event\_count = 200, stage\_counts = \{shared\_initial\_generation: 200\}, and model\_calls\_used = 200, confirming the enforced single shared invocation per task and that no separate condition‑specific sampling occurred.

Task bank, validators, and holdout emulator. The task bank is a curated, machine‑checkable intent\ensuremath{\rightarrow}configuration set (200 tasks) in which each task encodes initial and expected final states plus workflow policies (e.g., make\_before\_break\_uplink\_switchover, must\_be\_down\_for\_fields). The primary deterministic checker (deterministic\_netops\_validator\_v5, validator\_version = "5.0") runs syntactic parsing, intent/constraint enforcement, and structured diagnostics (fields available in validator traces include valid flag, violation\_count, violation\_codes, parsed\_command\_count, required\_command\_count, and difficulty annotations). A Mininet‑based data‑plane emulator was preregistered as a secondary disagreement detector applied to a randomized 10\% holdout (20 tasks) to detect validator--emulator mismatches; the executed contamination\_summary artifact reports zero disagreement flags for that holdout. The preregistration specified deterministic treatment rules: pairs flagged for validator--emulator disagreement would be excluded from the primary analysis per the contamination\_plan; sensitivity analyses for alternative treatments were preregistered but not invoked because no flags occurred.

Repair policy, budgeting, and missingness handling. Per the preregistration, when the primary validator reported violations the pipeline was permitted to issue up to one automated repair prompt to the same LLM for that task (maximum repair calls across the sample = 200). Repair prompts, if used, would be captured in repair\_provider\_trace fields in the episode artifacts; none appear in the archived episodes because no repairs were applied. Missingness and exclusion rules were deterministic: pairs where neither condition completed because of infrastructure failures would be excluded and logged; pairs with contamination flags (validator--emulator disagreement) would be excluded from the primary confirmatory test. The execution manifest and missingness\_summary show complete\_pair\_count = 200 and zero failed episodes, so the planned missingness handling was not invoked in practice.

Execution accounting and artifact capture. The autonomous run executed between 2026‑08‑30T11:50:34.530059Z and 2026‑08‑30T12:12:07.541650Z. The adapter enforced the preregistered maximum model calls (maximum\_model\_calls = 400) and per‑task caps (initial generation = 1 call; repair calls <= 1). The run consumed 200 model calls (one shared initial generation per task), with cache\_miss\_count = 200 and cache\_hit\_count = 0 as recorded in the deterministic reconciliation artifacts. All model calls, prompts, raw responses, validator traces, provider traces, and analysis artifacts (execution/raw\_results.jsonl, analysis/paired\_contingency\_table.csv, analysis/condition\_summary.csv, analysis/contamination\_summary.csv, analysis/analysis\_log.jsonl, execution/model\_configuration.json) were archived; representative per‑task artifacts (e.g., execution/scoring/task-000001-baseline.json) show validator\_trace fields and scoring outcomes. The analysis pipeline executed the exact McNemar test and bootstrap CI generation exactly as preregistered; analysis/analysis\_log.jsonl records bootstrap\_resamples = 10000 and bootstrap\_seed = 1729.

Design rationale and implications for internal validity. The shared\_initial paired design was chosen to isolate the pure marginal effect of validation+repair on a fixed generator output; this reduces variance attributable to stochastic sampling and makes the confirmatory test focused on whether repair can convert a validator‑failing candidate into a validator‑passing one for the same initial output. The trade‑off is explicit: the design sacrifices generality across multiple independent generator draws for stronger internal control. Preregistration, deterministic adapter semantics, and frozen stopping rules were used to avoid post‑hoc analytic choices. Deterministic reconciliation artifacts confirm pair\_accounting\_consistent = true and all\_deterministic\_consistency\_checks\_passed = true.

Limitations baked into the methodology. The methodology intentionally restricts scope: a single model, a deterministic prompt/sampling policy, a one‑repair budget, and a curated task bank. These constraints reflect an explicit confirmatory question about marginal GVR effect under a tightly controlled generator policy rather than an evaluation of alternative sampling, multi‑step repair strategies, or cross‑model generality. When interpreting outcomes, readers should account for these preregistered methodological choices documented in the archived plan.

\section{Results}
Execution outcomes. The run completed all planned episodes: completed\_episode\_count = 400 (shared initial generation reused across conditions) and complete\_pair\_count = 200. Adapter accounting shows model\_calls\_used = 200 (one shared initial generation per task) and cache\_miss\_count = 200. No repair calls were issued because every initial candidate passed deterministic\_netops\_validator\_v5. Condition summaries (analysis/condition\_summary.csv) show baseline successes = 200/200 and guarded successes = 200/200 (success\_rate = 1.0 for both). The paired contingency (analysis/paired\_contingency\_table.csv) is n11 = 200, n10 = 0, n01 = 0, n00 = 0.

Primary inferential analysis. With zero discordant pairs the observed paired difference estimate = 0.0. The exact McNemar two‑sided p = 1.0. The bootstrap percentile 95\% CI (10,000 resamples; seed = 1729) is [0.0, 0.0]. Under the preregistered estimand and test the prespecified >=50\% relative‑reduction hypothesis is logically untestable on this sample because there were no baseline validator failures to reduce.

Representative diagnostics. Representative per‑task artifacts (for example, pair‑000001) confirm that the shared initial generator outputs matched the task required\_sequence and that validator\_trace.validation\_after.valid = true with violation\_count = 0. Difficulty annotations in validator traces (examples: make\_before\_break\_uplink\_switchover, paired\_link\_atomic\_mtu\_change) and parsed\_command\_count/required\_command\_count show that tasks exercised medium‑to‑very\_hard workflow constraints despite generating validator‑satisfying candidates. The Mininet 10\% holdout (20 tasks) executed as preregistered; analysis/contamination\_summary.csv reports zero validator--emulator disagreements, and deterministic reconciliation artifacts report all consistency checks passed (pair\_accounting\_consistent = true).

Unexecuted branches. Planned sensitivity branches for flagged contamination treatments and alternative missingness handling were not invoked because no tasks were flagged and no episodes failed. Secondary exploratory estimands (average repair iterations, GVR‑introduced failure rate, model‑call cost per successful repair) are empty in the analysis bundle because no repair attempts occurred; the analysis artifacts record these secondary\_results as empty.

\section{Discussion}
Conditional interpretation and statistical mechanics. The confirmatory null (zero baseline validator failures) is an artifact‑grounded outcome: baseline\_success\_rate = 1.0 (200/200) and guarded\_success\_rate = 1.0 (200/200), producing a paired contingency with n11 = 200 and no discordant cells. In paired binary inference this configuration is degenerate for McNemar's test---the exact two‑sided p = 1.0 and the bootstrap percentile 95\% CI = [0.0, 0.0] (bootstrap parameters: resamples = 10,000, seed = 1729). Those numbers are not evidence that repair is useless generally; they are a precise statement that, on this curated workload and under the frozen shared‑initial generation policy, there were no validator‑detectable failures for the GVR pipeline to reduce. The key interpretive point is therefore conditional: the marginal utility of automated repair is a function of baseline validator failure prevalence and validator coverage for the targeted workload.

Artifact‑grounded diagnostics explaining the null. Several archived artifacts document why initial candidates passed at scale. First, each task encodes explicit required\_sequence and workflow constraints; the primary validator checks those constraints and reports parsed\_command\_count and required\_command\_count for each task (representative traces show required\_command\_count values matching parsed\_command\_count). Second, difficulty annotations in validator traces (labels such as "make\_before\_break\_uplink\_switchover", "paired\_link\_atomic\_mtu\_change") show that the task bank exercised nontrivial dependency patterns even when outputs were correct; correctness under these labels indicates the generator/prompt pair produced sequences that satisfied the machine‑checkable constraints. Third, adapter accounting (deterministic\_reconciliation) shows that hosted\_model\_call\_event\_count = 200 and stage\_counts = \{shared\_initial\_generation: 200\}, confirming a single deterministic generation per task was shared across conditions---this design choice removed sampling variability as a pathway to observe alternative (potentially failing) candidates. Fourth, the preregistered Mininet secondary check ran on a randomized 10\% holdout (20 tasks) and the contamination summary (analysis/contamination\_summary.csv) records zero disagreements between the primary validator and emulator; this artifact verifies that, in the holdout sample, validator verdicts matched the emulator's checks as executed in this run. We also disclose a reproducibility gap: the labeled artifact set does not include an explicit selection seed for the Mininet holdout, which prevents bit‑for‑bit reconstruction of the exact holdout selection even though the contamination summary reports zero disagreements.

Methodological tradeoffs: shared‑initial pairing versus unconditional failure measurement. The shared\_initial paired design was chosen to isolate the incremental effect of invoking validator+repair on the same generator output. That isolation gains internal control at the cost of observing only the conditional effect on a fixed sample: it cannot capture gains obtainable by re‑sampling the LLM (e.g., producing alternative candidates that a repair could act upon) or by employing nonzero temperature sampling to elicit diverse failure modes. For practitioners or evaluators whose goal is to estimate unconditional operational failure rates (the risk of an LLM pipeline producing an unsafe candidate on any application), a design with independent per‑condition sampling and multiple draws per task will produce better estimates of unconditional baseline failure prevalence and of the repair mechanism's ability to convert some failing draws into passes. Our artifact set supports both measurements in principle, but this confirmatory run intentionally prioritized a paired, shared‑initial comparison per the preregistered estimand.

Validator coverage and emulator corroboration. Deterministic validators are necessary but not sufficient oracles for operational safety: they codify a set of syntactic and policy checks (here, deterministic\_netops\_validator\_v5) that map to task‑bank constraints. The Mininet emulator was included as a pragmatic secondary detector to reveal validator blind spots that manifest at the data‑plane level; the zero disagreements observed in the 20‑task holdout demonstrate no such blind spots were detected in that sample for this run, but holdout scale and the missing labeled selection seed limit how strongly that result can be generalized. A robust operational assurance program should therefore combine diverse validators (vendor‑specific semantics, richer device models), larger emulation or testbed sampling, and adversarial/perturbation tests to probe validator blind spots; the artifacts here illustrate where such secondary checks succeeded (no disagreements) and where exact reproducibility of the holdout draw is currently unverifiable (missing seed).

Design implications for guardrail engineering and operational deployment. The neutral outcome suggests practitioners evaluate two complementary axes before investing heavily in automated repair: (1) measure baseline validator failure prevalence on representative workloads under the intended generation policy; (2) assess validator completeness via emulation or selective real‑device checks. If baseline failures are already rare under the production generator/prompt policy, then (a) a single‑pass repair budget delivers limited marginal safety returns for validator‑detected failures, and (b) resources may be better allocated to expanding validator fidelity, adversarial testing, or introducing multi‑sample generation and model ensembles to surface latent failure modes. Conversely, when baseline failures are nontrivial, GVR with an appropriately sized repair budget and well‑designed repair prompts can reduce validator‑detected errors; the archived preregistration captures how repair budgeting and maximum repair calls were constrained here (up to one repair call per task), and the absence of repair events means the exploratory statistics for repair iterations and model‑call costs remained empty in the analysis artifacts.

Threats to validity and concrete artifact links. Internal validity is strengthened by preregistration, deterministic adapter enforcement, and deterministic reconciliation artifacts that confirm episode accounting (complete\_pair\_count = 200, model\_calls\_used = 200). Construct validity remains conditional on validator expressiveness: representative validator\_trace fields (validator\_id, validator\_version, violation\_count, parsed\_command\_count) show what the validator measured, but validators do not represent all device idiosyncrasies. External validity is limited by the single hosted model (gpt‑5‑mini), the shared\_initial deterministic sampling policy, and the curated task bank; these constraints are recorded in execution/model\_configuration.json and the preregistration artifact. Conclusion validity: with zero discordant observations the primary confirmatory test is uninformative about the preregistered 50\% relative reduction; the observed numbers are nonetheless exact and reproducible within the archived bundle (aside from the Mininet‑holdout seed gap noted above).

Practical recommendations for future confirmatory studies. To measure GVR effectiveness where baseline failure prevalence is low, future preregistered runs should (i) adopt a mixed design that includes independent per‑condition multi‑sample draws to estimate unconditional failure rates; (ii) increase the repair budget and allow multi‑step repairs or human‑in‑the‑loop escalation to exercise more repair pathways; (iii) scale secondary emulation checks and record selection seeds for holdouts to enable exact reproducibility; and (iv) evaluate multiple models or ensembles to probe model‑dependent brittleness. The archived artifacts and the deterministic reconciliation logs provide a template for implementing these changes while preserving auditability.

\section{Limitations}
\begin{itemize}
\item Task‑bank scope: the confirmatory sample used a curated synthetic/public intent\ensuremath{\rightarrow}configuration bank (n = 200); results therefore reflect this bank and may not generalize to broader production workloads or adversarial distributions.
\item Single‑model and shared‑initial setting: only gpt‑5‑mini (hosted) with adapter‑enforced shared‑initial generation was tested (execution artifacts record hosted\_model\_call\_event\_count = 200 and stage\_counts shared\_initial\_generation = 200). Outcomes may differ across model families, independent per‑condition sampling, nonzero temperature sampling, or larger repair budgets.
\item Validator coverage and oracle limitations: the primary deterministic validator enforces a defined set of syntax/intent constraints; validators can fail to capture real device idiosyncrasies, vendor‑specific behaviors, or emergent failure modes detectable only by end‑to‑end deployment.
\item Adversarial robustness untested: this confirmatory run did not include active adversarial injection or red‑team tool‑description attacks; prior audits indicate such vectors can bypass naive defenses and remain a separate concern.
\item Reproducibility gap for contamination check: the randomized holdout selection seed for the Mininet contamination check is absent from the labeled artifacts, preventing exact reconstruction of the holdout selection.
\item Single repair budget: the GVR pipeline allowed at most one automated repair prompt per task; multi‑step repair strategies, different repair budgets, or human‑in‑the‑loop approvals were not exercised here.
\end{itemize}

\section{Conclusion}
We preregistered and executed a sealed paired study comparing LLM‑only generation to a deterministic GVR pipeline on 200 NetOps intent\ensuremath{\rightarrow}configuration tasks using gpt‑5‑mini and deterministic\_netops\_validator\_v5. Both pipelines produced validator‑passing configurations for every task (n11 = 200; no discordant pairs), resulting in an observed paired difference of 0.0, exact McNemar p = 1.0, and bootstrap 95\% CI = [0.0, 0.0]. The preregistered >=50\% relative‑reduction hypothesis could not be evaluated because baseline validator failures were absent. The artifact‑grounded outcome clarifies that GVR's marginal value is workload dependent: when baseline validator‑detectable failures are rare, automated repair yields no measured benefit for those validator‑detected errors. Future studies should target workloads with nontrivial baseline failure rates, expand validator fidelity with richer emulation or real‑device checks, evaluate multi‑sample and multi‑model policies, and explore multi‑step or human‑in‑the‑loop repair budgets to characterize practical operational gains. All execution artifacts and analysis outputs are archived for independent inspection under the preregistration id cited above.

\section*{Disclosure Statement}
Disclosure (mandatory). This study executed under the autonomous post‑lock preregistration cnsms2026‑gvr‑effectiveness‑02. Declared model: gpt‑5‑mini (provider label: openai\_responses) invoked via the hosted\_netops\_gvr\_v1 adapter. The exact initial master prompt is archived at provenance/master\_prompt.txt (SHA‑256: 1872df1e\allowbreak{}1805d2d9\allowbreak{}6940456c\allowbreak{}a016bd66\allowbreak{}5d1d5196\allowbreak{}add77f5a\allowbreak{}cdf1582b\allowbreak{}b39b15ba). Topic constraints: the human Research Director limited the study to Generative AI and LLMs for NetOps focusing on safety/reliability of generated network changes. Frameworks and tools: orchestration used hosted\_netops\_gvr\_v1 and the deterministic task generator deterministic\_netops\_task\_generator\_v5. Primary deterministic validator: deterministic\_netops\_validator\_v5 (intent/constraint checks). A Mininet‑based emulator was preregistered and executed as a secondary disagreement detector on a randomized 10\% holdout (20 tasks); the artifact analysis/contamination\_summary.csv shows zero disagreements for that holdout. Preregistration and freeze: the experimental plan (estimand, sample size n = 200, paired design, stopping rules, max model calls = 400, repair budget = 1 per task, shared‑initial generation semantics) was frozen before any model calls. Autonomous execution and artifacts: the run executed autonomously under the frozen plan (start: 2026‑08‑30T11:50:34.530059Z; end: 2026‑08‑30T12:12:07.541650Z); model calls, prompts, seeds, validator traces, emulator traces, logs, and the artifact manifest are archived in the execution bundle. Model‑call budget and usage: 200 model calls were used (one shared initial generation per task); up to 200 repair calls were allowed but none were applied because initial candidates passed the validator. Human involvement: the human Research Director selected the topic family and pre‑lock constraints; no human manual editing or selective rewriting of technical manuscript content occurred after the autonomy lock. Deviations and restarts: no post‑preregistration design changes or technical restarts occurred. Known reproducibility gap: the explicit randomized selection seed used to draw the Mininet 10\% holdout is not present as a labeled artifact in the archive; this absence prevents exact deterministic reconstruction of the drawn holdout despite the contamination summary reporting no disagreements. The master prompt archive reference above is the immutable prompt required by the track.

\begin{thebibliography}{99}
\bibitem{ref1} Wayne Xin Zhao, Kun Zhou, Junyi Li, Tianyi Tang, Zican Dong, Yupeng Hou, Beichen Zhang, Yingqian Min, Junjie Zhang, Peiyu Liu, Xiaolei Wang, Yifan Du, Yushuo Chen, Yushuo Chen, Zhipeng Chen, Jinhao Jiang, Ruiyang Ren, Yifan Li, Xinyu Tang, Peiyu Liu, Yiwen Hu, Jian‐Yun Nie, Ji-Rong Wen, "A Survey of Large Language Models", 2026, 10.1007/s11704-026-60308-3
\bibitem{ref2} Changjie Wang, Mariano Scazzariello, Alireza Farshin, Simone Ferlin, Dejan Kostić, Marco Chiesa, "NetConfEval: Can LLMs Facilitate Network Configuration?", 2024, 10.1145/3656296
\bibitem{ref3} Oscar G. Lira, Oscar Maurício Caicedo Rendón, Nelson L. S. da Fonseca, "Large Language Models for Zero Touch Network Configuration Management", 2024, 10.1109/mcom.001.2400368
\bibitem{ref4} Chirag Shah, Ryen W. White, Reid Andersen, Georg Buscher, Scott Counts, Sarkar Snigdha Sarathi Das, Ali Montazeralghaem, Sathish Manivannan, J. Neville, Nagu Rangan, Tara Safavi, Siddharth Suri, Mengting Wan, Leijie Wang, Longqi Yang, "Using Large Language Models to Generate, Validate, and Apply User Intent Taxonomies", 2025, 10.1145/3732294
\bibitem{ref5} omar altawil, Dr. Mustafa Al-Fayoumi, "Configuration-Level Semantic Brittleness in Tool-Using LLM Agents: A Factor-Ranking Study of Persona, History, and Tool Definition", 2026, 10.2139/ssrn.7203631
\bibitem{ref6} Mohammadreza Rashidi, "Measuring How LLM Tool Descriptions, Cross-Tool Ambiguity, and Action Chains Compose into Excessive Agency Exploits", 2026, 10.21203/rs.3.rs-10646209/v1
\bibitem{ref7} Yoshihiro Ando, "Secure-by-Default Guardrails for MCP-Based Tool Use in Multi-Modal LLM Agents", 2026, 10.36227/techrxiv.177006109.97957916/v1
\bibitem{ref8} Sebastian Farquhar, Jannik Kossen, Lorenz Kuhn, Yarin Gal, "Detecting hallucinations in large language models using semantic entropy", 2024, 10.1038/s41586-024-07421-0
\bibitem{ref9} Shuyin Ouyang, Jie M. Zhang, Mark Harman, Meng Wang, "An Empirical Study of the Non-Determinism of ChatGPT in Code Generation", 2024, 10.1145/3697010
\bibitem{ref10} Oghenekeno Hilkiah Okula, ""The Era of AI Agents for Network Engineering": Intent-Aware Auto Remediation for Network Disruption or Failures Using AI Agents, Large Language Models (LLM) and Model Context Protocol (MCP) Mechanisms", 2026, 10.36227/techrxiv.177004277.71795055/v1
\end{thebibliography}

\end{document}
